% Image credits for the photographs and AI illustrations of Book 1 (Primary and Middle School), French edition.
% Machine-readable license records live in images/book1/CREDITS.md; keep this
% file in sync with frontmatter/image-credits.tex (one \item per item there).
\chapter*{Crédits des images}

Les dessins, cartes et graphiques de ce livre sont originaux~; les cartes
et graphiques tracés à partir de données publiées en citent la source dans
la légende ou dans cette liste. Les photographies sont utilisées sous leurs
licences libres respectives, avec nos remerciements à leurs auteurs~; les
liens complets vers les sources et les adresses des licences de chaque
photographie figurent dans \texttt{images/book1/CREDITS.md}, dans le dépôt
du livre.

\bigskip
À côté des photographies et des dessins originaux, certaines figures sont
des illustrations générées par IA, produites avec le générateur d'images
d'OpenAI à partir de consignes rédigées par les éditeurs du livre, et
vérifiées pour leur exactitude scientifique avant d'être retenues. Elles
ne montrent que des scènes génériques, jamais un lieu réel, une personne
réelle ou un spécimen réel. La consigne complète de chaque image générée
est consignée dans \texttt{images/book1/ai/PROMPTS.md}, dans le dépôt.

\bigskip
\noindent Photographies et données, par chapitre~:
\begin{itemize}\small
  \item Chap.~\ref{ch:g2:water-on-earth} (l'eau sur la Terre)~: la Terre vue d'Apollo~17, 7~décembre 1972 («~The Blue Marble~»)~; NASA, domaine public (via Wikimedia Commons). Les 100 carrés sont calculés d'après le relief NOAA ETOPO1~; les 100 seaux d'après le tableau de la Water Science School de l'U.S. Geological Survey (d'après Shiklomanov, 1993).
  \item Chap.~\ref{ch:g3:stones-pebbles-sand} (cailloux, galets et sable)~: un mur de pierres sèches en granite~; photographie de Tony Atkin, Geograph, CC BY-SA 2.0 (via Wikimedia Commons). Des toits en dalles de pierre en Suisse~; photographie de Michael C. Rygel, CC BY-SA 3.0 (via Wikimedia Commons).
  \item Chap.~\ref{ch:g3:continents-and-oceans} (le globe)~: les deux globes et le planisphère sont dessinés d'après les côtes de Natural Earth (domaine public).
  \item Chap.~\ref{ch:g4:clouds-rain-snow} (nuages, pluie et neige)~: des cumulus au-dessus du Kansas (S.~Corfidi), des plaques de cirrus (R.~F.~Kresge), une couche de stratus autour de sommets montagneux, un orage au-dessus du désert de l'Arizona (collection de B.~Hooke), un très gros grêlon~; tous NOAA Photo Library, domaine public (via Wikimedia Commons).
  \item Chap.~\ref{ch:g4:volcanoes-earthquakes-first-look} (volcans et séismes, un premier regard)~: de la lave entrant dans l'océan au Kilauea, USGS Hawaiian Volcano Observatory~; le mont St.~Helens, 18~mai 1980, photographie d'Austin Post, USGS~; une maison arrachée à ses fondations par le séisme de Loma Prieta de 1989, photographie de J.~K.~Nakata, USGS~; tous domaine public (via Wikimedia Commons). La carte utilise la liste des volcans holocènes du Global Volcanism Program de la Smithsonian Institution et le catalogue des séismes de l'U.S. Geological Survey.
  \item Chap.~\ref{ch:g4:world-physical-map} (la carte physique du monde)~: l'Himalaya et l'Amazone photographiés depuis la Station spatiale internationale, NASA, domaine public (via Wikimedia Commons). La carte physique est dessinée d'après le relief NOAA ETOPO1, avec les fleuves, les lacs et les côtes de Natural Earth (domaine public).
  \item Chap.~\ref{ch:g5:water-shapes-the-land} (l'eau façonne les paysages)~: le Mamoré photographié depuis la Station spatiale internationale (ISS007-E-10797) et le delta du Nil vu par l'instrument MODIS du satellite Terra, NASA, domaine public (via Wikimedia Commons).
  \item Chap.~\ref{ch:g5:weather-and-climate} (le temps et le climat)~: le désert de Sonora, U.S. National Park Service, Saguaro National Park, domaine public (via Wikimedia Commons). Les diagrammes climatiques sont tracés d'après les U.S. Climate Normals 1991--2020 du NOAA NCEI (domaine public).
  \item Chap.~\ref{ch:g5:natural-hazards} (les aléas naturels)~: l'œil de l'ouragan Isabel photographié depuis la Station spatiale internationale par Ed Lu, NASA, domaine public~; le glissement de terrain de 2014 dans l'État de Washington, Mark Reid, U.S. Geological Survey, domaine public (tous deux via Wikimedia Commons).
  \item Chap.~\ref{ch:g6:earth-a-planet} (la Terre, une planète)~: Vénus vue par Mariner~10 (NASA/JPL-Caltech)~; la Terre vue d'Apollo~17 (NASA)~; \emph{Earthrise} vu d'Apollo~8 (NASA, Bill Anders)~; Mars et Jupiter vues par le télescope spatial Hubble (NASA, ESA, the Hubble Heritage Team et A.~Simon)~; tous domaine public (via Wikimedia Commons). Le \emph{Pale Blue Dot} vu par Voyager~1 en 1990, version de 2020 (NASA/JPL-Caltech, PIA23645, domaine public). Données des planètes d'après la fiche planétaire de la NASA.
  \item Chap.~\ref{ch:g6:latitude-longitude} (latitude et longitude)~: le chronomètre H4 de Harrison au National Maritime Museum de Greenwich, photographie de Colonel Warden, CC BY-SA 3.0 (via Wikimedia Commons). Côtes et villes de Natural Earth (domaine public).
  \item Chap.~\ref{ch:g6:rocks-and-minerals} (roches et minéraux)~: granite, basalte, calcaire fossilifère, grès, marbre, argilite, un cristal de quartz et une feuille de mica muscovite, photographies de James St.~John, CC BY 2.0~; un toit d'ardoises, photographie d'Øyvind Holmstad, CC0 (tous via Wikimedia Commons).
  \item Chap.~\ref{ch:g6:inside-the-earth} (l'intérieur de la Terre)~: Inge Lehmann en 1932, photographie d'Even Neuhaus, Bibliothèque royale du Danemark, domaine public (via Wikimedia Commons). Les couches, les masses volumiques et les vitesses des ondes sont tirées du Preliminary Reference Earth Model (A.~M. Dziewonski et D.~L. Anderson, 1981).
  \item Chap.~\ref{ch:g6:water-underground} (l'eau souterraine)~: Big Spring, Missouri, photographie de Johnston9494, CC BY-SA 4.0 (via Wikimedia Commons). Porosités d'après R.~C. Heath, \emph{Basic Ground-Water Hydrology}, U.S. Geological Survey (1983)~; parts d'eau douce d'après la Water Science School de l'USGS.
  \item Chap.~\ref{ch:g7:plate-tectonics} (la tectonique des plaques)~: les continents du Sud réassemblés, avec les bandes de fossiles de \emph{Cynognathus}, \emph{Lystrosaurus}, \emph{Mesosaurus} et \emph{Glossopteris}~; U.S. Geological Survey, \emph{This Dynamic Earth}, domaine public.
  \item Chap.~\ref{ch:g7:plate-tectonics}~: la faille de San Andreas à travers la Carrizo Plain~; photographie de Robert E. Wallace, U.S. Geological Survey, domaine public.
  \item Chap.~\ref{ch:g7:plate-tectonics}~: limites des plaques d'après P.~Bird (2003), GeoJSON de H.~Ahlenius (ODC-By 1.0)~; côtes de Natural Earth (domaine public)~; séismes du catalogue de l'U.S. Geological Survey~; volcans du Global Volcanism Program de la Smithsonian Institution.
  \item Chap.~\ref{ch:g7:volcanoes} (les volcans)~: le mont St.~Helens vu de Johnston Ridge le 17~mai et le 10~septembre 1980, photographies de Harry Glicken, U.S. Geological Survey~; une fontaine de lave au Kilauea, 1983, et le Mauna Loa vu d'avion, 1985, photographies de J.~D. Griggs, U.S. Geological Survey~; la colonne éruptive du mont Pinatubo, 12~juin 1991, photographie de Dave Harlow, U.S. Geological Survey~; tous domaine public (via Wikimedia Commons). La carte et les décomptes utilisent la liste des volcans holocènes du Global Volcanism Program de la Smithsonian Institution~; limites des plaques d'après P.~Bird (2003).
  \item Chap.~\ref{ch:g7:earthquakes} (les séismes)~: un village près de la côte de Sumatra après le tsunami du 26~décembre 2004, photographie de l'U.S. Navy par Philip A. McDaniel~; les escarpements de faille du séisme de Borah Peak de 1983, photographie de R.~E. Wallace, U.S. Geological Survey~; la route affaissée dans le lac Hebgen, 1959, photographie d'I.~J. Witkind, U.S. Geological Survey~; tous domaine public (via Wikimedia Commons). Les décomptes et les profondeurs utilisent le catalogue des séismes de l'U.S. Geological Survey~; le tableau des intensités est abrégé de l'échelle de Mercalli modifiée de l'U.S. Geological Survey.
  \item Chap.~\ref{ch:g7:sediments-and-fossils} (sédiments et fossiles)~: le Grand Canyon vu de Hermits Rest, photographie NPS de Kristen M. Caldon, domaine public~; une ammonite, un trilobite et des rides de courant dans le parc national de Capitol Reef, photographies de James St.~John, CC BY 2.0~; des empreintes de dinosaures à Dinosaur Ridge, photographie de Dexter Perkins, CC0 (tous via Wikimedia Commons).
  \item Chap.~\ref{ch:g7:geological-time} (les âges de la Terre)~: James Hutton, portrait par Henry Raeburn, domaine public (via Wikimedia Commons). Les âges et les couleurs de l'échelle des temps sont ceux de la Charte chronostratigraphique internationale (Commission internationale de stratigraphie).
  \item Chap.~\ref{ch:g8:the-atmosphere} (l'atmosphère)~: le lâcher d'un ballon-sonde à Topeka, au Kansas, photographie de Stephen Corfidi, NOAA/National Weather Service~; le bord de la Terre au coucher du Soleil vu de la Station spatiale internationale, NASA, ISS Expedition~23~; tous deux domaine public (via Wikimedia Commons). Les profils sont calculés d'après l'\emph{U.S. Standard Atmosphere, 1976} (NOAA, NASA et U.S. Air Force).
  \item Chap.~\ref{ch:g8:winds-global-circulation} (les vents et la circulation générale)~: les Amériques vues par le satellite GOES-16 le 26~juillet 2022, NOAA, domaine public (via Wikimedia Commons). Les côtes de la carte des alizés viennent de Natural Earth (domaine public).
  \item Chap.~\ref{ch:g8:fronts-and-weather-maps} (masses d'air, fronts et cartes météorologiques)~: une dépression au large de l'Islande vue par le satellite Terra, NASA/GSFC, MODIS Rapid Response Team, Jacques Descloitres, domaine public (via Wikimedia Commons). Les côtes de la carte des masses d'air viennent de Natural Earth (domaine public).
  \item Chap.~\ref{ch:g8:the-oceans} (les océans)~: le HMS \emph{Challenger} se préparant à sonder, 1872, gravure tirée des rapports de l'expédition du Challenger reproduite dans W.~A. Herdman, \emph{Founders of Oceanography} (1923)~; le Gulf Stream vu par l'instrument MODIS du satellite Aqua, NASA Earth Observatory~; tous deux domaine public (via Wikimedia Commons). Les profils du fond océanique sont tracés d'après le relief NOAA ETOPO1~; les côtes des cartes des gyres et des bouées viennent de Natural Earth (domaine public).
  \item Chap.~\ref{ch:g8:ice-on-earth} (la glace sur la Terre)~: le glacier Pedersen en 1990 et en 2016, photographies de M.~Tetreau et D.~Kurtz, U.S. National Park Service~; le bord d'une plate-forme de glace au large de la péninsule Antarctique, NASA, Operation IceBridge (M.-J. Viñas)~; tous domaine public (via Wikimedia Commons). Le graphique de la banquise utilise le NSIDC Sea Ice Index, version~4 (F.~Fetterer et al., 2025, doi:10.7265/a98x-0f50).
  \item Chap.~\ref{ch:g9:climate-past-and-present} (le climat de la Terre, hier et aujourd'hui)~: un tronçon de la carotte de glace GISP2 du Groenland, NOAA Paleoclimatology, domaine public (via Wikimedia Commons). Les séries sont tracées d'après NASA GISTEMP version~4, la série du dioxyde de carbone de Mauna Loa (NOAA Global Monitoring Laboratory et Scripps Institution of Oceanography), la température d'EPICA Dôme~C (J.~Jouzel et al., 2007), la série composite du dioxyde de carbone des carottes antarctiques (B.~Bereiter et al., 2015) et le niveau de la mer du NOAA Laboratory for Satellite Altimetry~; les affirmations de l'évaluation sont celles du sixième rapport d'évaluation du GIEC, groupe de travail~I, résumé à l'intention des décideurs (2021).
  \item Chap.~\ref{ch:g9:natural-risks} (vivre avec les risques naturels)~: la carte des aléas du Nevado del Ruiz, d'après T.~L. Wright et T.~C. Pierson, \emph{Living with Volcanoes}, U.S. Geological Survey Circular 1073 (1992)~; une bouée d'alerte aux tsunamis inspectée par l'équipage de l'USCGC \emph{William Hart}, U.S. Coast Guard, 2021~; toutes deux domaine public (via Wikimedia Commons).
  \item Chap.~\ref{ch:g9:earth-resources} (les ressources de la Terre)~: la mine de Chuquicamata vue par ASTER, NASA/\allowbreak METI/\allowbreak AIST/\allowbreak Japan Space Systems et U.S./Japan ASTER Science Team, 2007, domaine public~; une veine de charbon dans un talus routier du comté de Belmont, Ohio, photographie de James St.~John, CC BY 2.0 (toutes deux via Wikimedia Commons). Les facteurs de concentration viennent de D.~Perkins et al., \emph{Mineralogy}, et de S.~Earle, \emph{Environmental Geology}.
  \item Chap.~\ref{ch:g9:landforms} (les formes du relief)~: les Seven Sisters, photographie d'Arild Vågen, 2017, CC BY-SA 4.0~; des dunes en croissant sur la côte du sud du Brésil, photographie prise par un astronaute de la NASA depuis la Station spatiale internationale, 2017, domaine public~; des stalactites à Carlsbad Cavern, photographie du National Park Service par Daniel Leifheit, 2018, domaine public (toutes via Wikimedia Commons).
  \item Chap.~\ref{ch:g9:maps-of-the-earth} (cartographier la Terre entière)~: le delta de la Lena, image Landsat de 2000, USGS EROS Data Center via NASA Visible Earth, domaine public~; un satellite GPS, vue d'artiste de la NASA, domaine public (tous deux via Wikimedia Commons). Les côtes des planisphères viennent de Natural Earth (domaine public).
\end{itemize}
